\section{Introduction}
\label{sec:intro}

Software rarely exists in a single, immutable form. The same functionality is routinely re-expressed across languages, libraries, and coding styles as systems are ported, hardened, and modernized. This plurality is not incidental: it is a defensive and engineering asset. \emph{Software diversity} raises the cost of exploitation by denying an attacker a single, uniform target~\cite{larsen2014sok}, and dependency modernization, replacing an outdated or vulnerable component with a safer or faster one, is a recurring maintenance burden in long-lived deployments. Nowhere is this burden more acute than in \emph{critical-infrastructure} software, such as the code that estimates measurements on a factory floor or drives an industrial control system (ICS), where a component must be replaced or diversified without altering the process it governs. In such settings the central requirement is exacting: the new program must behave \emph{exactly} like the old one, differing only where we intend it to differ.

For over two decades, the research community has studied the recognition side of this equivalence question under the banner of \emph{code clone detection}~\cite{roy2009comparison,sajnani2016sourcerercc,bin_sim_survey}, the task of deciding whether two programs are semantically equivalent despite syntactic differences. Detection underpins vulnerability search, malware triage, plagiarism analysis, and intellectual-property protection, and recent work, including our own prior study of Large-Language-Model (LLM)-assisted binary clone detection~\cite{agadakos2025bcd}, has shown that modern models, when augmented with program-analysis tooling, can reason about semantic equivalence with useful accuracy. The inverse problem, by contrast, has received little attention: rather than \emph{recognizing} that two programs are clones, can we \emph{construct} a clone of a given program on demand?

We argue that automated clone \emph{generation} is both timely and tractable. It is timely because LLMs have demonstrated strong performance across code synthesis, translation, and repair~\cite{chen2021codex,li2022alphacode,roziere2020transcoder}, and because agentic architectures now let these models act, observe the results of their actions, and iteratively refine their output~\cite{yao2023react,shinn2023reflexion}. Prior work has also shown that grounding LLMs in execution evidence improves the reasoning of such systems, letting them solve problems they would otherwise fail~\cite{agadakos2025bcd}.

It is tractable because generation, unlike detection, admits a \emph{constructive} correctness check: a candidate clone can simply be executed and compared against the original. The key insight is that an LLM need not produce a correct clone in one shot; it needs only to produce a candidate that a verification loop can execute, test, and repair until equivalence holds. Importantly, this shifts the burden from the model's fallible one-shot reasoning, a well-documented source of hallucination, onto concrete runtime evidence.

To this end, we present \sys, an agentic system that synthesizes a functional code clone from a given source program and a set of transformation requirements. \sys decomposes the problem into four cooperating, LLM-driven agents, each responsible for a distinct class of change and each guarded by its own execute-and-repair loop, with a fifth agent standing in for two of the four when the requirement calls for removing a dependency rather than substituting one. A Library-Substitution Agent replaces third-party dependencies (e.g., a cryptographic or numerical library) with functionally comparable alternatives that are more secure or faster, grounded in benchmark and vulnerability evidence rather than the model's own prior beliefs; a Glue-Synthesis Agent writes the adapter code that reconciles the program with the substituted dependency's interface; a Logic-Mutation Agent rewrites \emph{first-party} application logic, not third-party code, into behaviorally equivalent but textually divergent form; and a Verification and Consistency Agent gates every change against a deterministic completeness check and an independent model review, confirming the migration is genuine or routing each detected regression back to the agent responsible for it.

We evaluate \sys on software drawn from the industrial-measurement and infrastructure-control domain, where behavioral fidelity is a hard constraint rather than a convenience. Our contributions are:

\begin{itemize}
    \item We formulate \emph{automated code clone generation} as the constructive inverse of clone detection, and we identify verification-driven, agentic mutation as a practical means of achieving it.
    \item We propose \sys, a four-agent core architecture, with a fifth agent for programs that have no library to substitute, that separates dependency substitution, glue synthesis, first-party logic mutation, and cross-cutting equivalence verification, with a two-tier (local execute-and-repair and global oracle-based) correctness discipline.
    \item We develop a prototype whose global correctness signal is a five-layer source-level equivalence oracle: import validation, recorded golden vectors, program invariants, the subject's own test suite, and direct execution of the substituted factory path against the real library. Extending this to a binary-level oracle via our prior clone detector~\cite{agadakos2025bcd} is a direct next step we outline in Section~\ref{sec:impl}, not something the results reported here depend on.
    \item We evaluate \sys against three baselines on industrial measurement and control software and report migration success, an independent blind review of clone fidelity, verification cost, and structural divergence.
\end{itemize}

The rest of this paper is organized as follows. Section~\ref{sec:background} provides background on code clones, agentic LLMs, and the industrial-control setting. Section~\ref{sec:approach} presents the \sys architecture and its two verification tiers. Section~\ref{sec:impl} describes our prototype. Section~\ref{sec:results} reports our evaluation. Section~\ref{sec:limitations} discusses \sys's limitations and scope. Section~\ref{sec:related} surveys related work, and Section~\ref{sec:conclusions} concludes.
